\clearpage
\VerdierPhysicalPageAnchor{0238}%
\begingroup\small\setlength{\parskip}{0pt}\setlength{\emergencystretch}{1.5em}
\setlength{\abovedisplayskip}{4pt plus 2pt minus 2pt}%
\setlength{\belowdisplayskip}{4pt plus 2pt minus 2pt}%

\VerdierSemanticAnchor{ch4d-4-3-4}
\noindent\textbf{4.3.4.}\quad
$D(\mathcal A)$ から $D(\mathcal A)$ のスペクトル対象の圏への関手も、
用語を濫用して\emph{第二標準スペクトル対象関手}と呼ぶ。

$X$ を $\mathcal A$ の複体とし、
$\bigl(X(p,q),\delta(p,q,r)\bigr)$ を $X$ に付随する第二スペクトル対象とする。
(4.3.2) と (1.2.10) から、$D(\mathcal A)$ における標準同型を得る：

\VerdierSemanticAnchor{ch4d-4-3-4-1}
\begin{equation}
X(p,p+1)\xrightarrow{\sim}H^p(X)[-p].
\tag{4.3.4.1}
\end{equation}

\VerdierSemanticAnchor{ch4d-4-3-4-boundedness}
さらに、$X$ が $D^*(\mathcal A)$ の対象
（$*=-,+,b,$ または「空白」）であるとき、$X$ に付随する第二スペクトル対象は
実際には $D^*(\mathcal A)$ に値をとる。

\VerdierSemanticAnchor{ch4d-4-4}
\medskip\noindent\textbf{4.4. 通常のスペクトル系列。}

\VerdierSemanticAnchor{ch4d-4-4-1}
\noindent\textbf{4.4.1.}\quad
$\mathcal A$ を、必ずしもアーベル圏とは限らない加法圏とし、
$F\colon K(\mathcal A)\longrightarrow\mathcal B$ をアーベル圏
$\mathcal B$ へのコホモロジー的関手とする。また
$F^p\colon K(\mathcal A)\longrightarrow\mathcal B$ を、$F$ と
$K(\mathcal A)$ のシフト関手との合成から得られる関手
（$F^p(X)=F(X[p])$）とする。さらに $Y$ を $\mathcal A$ の複体とする。
関手 $F$ は、$Y$ に付随する $K(\mathcal A)$ に値をとる第一スペクトル対象
(4.1.4) を、$\mathcal B$ に値をとるスペクトル対象に移す。これから
（第II章 4.3.3）、$F$ と $Y$ の双方について関手的な（後者が
$\operatorname{comp}(\mathcal A)$ の中を動くとき）、$\mathcal B$ に値をとる
スペクトル系列を得る。これを\emph{関手 $F$ の複体 $Y$ に関する第一
スペクトル系列}という。このスペクトル系列には常に（第II章 4.3.3.1）で
導入した添字付けを用いる。

\VerdierSemanticAnchor{ch4d-4-4-2}
\medskip\noindent\textbf{命題 4.4.2.}\quad

\VerdierSemanticAnchor{ch4d-4-4-2-a}
\noindent\textit{a) 関手 $F$ の複体 $Y$ に関する第一スペクトル系列の項
$\mathrm I_1^{p,q}$ は次で与えられる：}
\[
\mathrm I_1^{p,q}=F^q(Y^p)
\]
\noindent\textit{（ここで $Y^p$ は、唯一の非零成分が次数零にあり、それが
複体 $Y$ の第 $p$ 成分である複体を表す）。対応する微分は}

\VerdierSemanticAnchor{ch4d-4-4-2-a-differential}
\[
\begin{aligned}
d_1&\colon\mathrm I_1^{p,q}\longrightarrow\mathrm I_1^{p+1,q},\\
d_1&=F^q(-d_Y^p)\colon F^q(Y^p)\longrightarrow F^q(Y^{p+1}).
\end{aligned}
\]

\VerdierSemanticAnchor{ch4d-4-4-2-b}
\noindent\textit{b) したがって、同じスペクトル系列の項
$\mathrm I_2^{p,q}$ は次のように記述できる：
$F^{q,\bullet}\colon\operatorname{comp}(K(\mathcal A))
\longrightarrow\operatorname{comp}(\mathcal B)$ を $F^q$ の複体への拡張、
$H_{\mathcal B}^p\colon\operatorname{comp}(\mathcal B)
\longrightarrow\mathcal B$ を第 $p$ コホモロジー対象関手、$Y^\bullet$ を
$\operatorname{comp}(K(\mathcal A))$ の次の対象とする：}

\VerdierSemanticAnchor{ch4d-4-4-2-b-complex}
\[
Y^\bullet=\cdots\longrightarrow Y^{p-1}
\xrightarrow{d_Y^{p-1}}Y^p\xrightarrow{d_Y^p}Y^{p+1}
\longrightarrow\cdots.
\]
\endgroup
